\documentclass[11pt,letterpaper]{article}
\usepackage[margin=0.72in,headheight=15pt,headsep=18pt,footskip=25pt]{geometry}
\usepackage[T1]{fontenc}
\usepackage{amsmath}
% Font settings from ../dgstyle.tex; retain this report's author-free layout.
\usepackage{newpxtext}
\usepackage[slantedGreek]{newpxmath}
\renewcommand{\ttdefault}{lmtt}
\linespread{1.06}
\usepackage{graphicx,xcolor,booktabs,array,caption,fancyhdr,enumitem}
\usepackage{microtype}
\usepackage[hidelinks]{hyperref}
\definecolor{ink}{HTML}{25343D}
\definecolor{teal}{HTML}{176F78}
\definecolor{gold}{HTML}{A16623}
\definecolor{muted}{HTML}{65767D}
\hypersetup{pdftitle={Helike as a Probe of a Quantum Yang-Mills Model},pdfauthor={},pdfsubject={Verified finite-cutoff GPU survey; convergence unresolved}}
\pagestyle{fancy}\fancyhf{}
\fancyhead[L]{\sffamily\scriptsize\color{muted} GIGI / COMPUTATIONAL RESEARCH}
\fancyhead[R]{\sffamily\scriptsize\color{muted} 23 NOVEMBER 2026}
\fancyfoot[L]{\sffamily\scriptsize\color{muted} HELIKE QUANTUM PROBE / FINITE-CUTOFF REPORT}
\fancyfoot[R]{\sffamily\small\thepage}
\renewcommand{\headrulewidth}{0.3pt}
\setlength{\parindent}{0pt}\setlength{\parskip}{7pt}
\setlist[itemize]{leftmargin=15pt,itemsep=3pt,topsep=4pt}
\captionsetup{font=small,labelfont=bf,justification=raggedright,singlelinecheck=false,skip=6pt,hypcap=false}
\newcommand{\eyebrow}[1]{{\sffamily\small\color{teal}#1}\par}
\newcommand{\pagetitle}[2]{\eyebrow{#1}\vspace{3pt}{\LARGE #2}\par\vspace{5pt}}
\newcommand{\fig}[2]{\begin{center}\includegraphics[width=\linewidth]{figures/#1.pdf}\captionof{figure}{#2}\end{center}}
\newcommand{\Htwo}{\mathrm{H2}}\newcommand{\Hzero}{\mathrm{H0}}
\newcommand{\runid}{\texttt{20261124T021903Z-c326a3c5}}
\begin{document}
\color{ink}
\eyebrow{EXPERIMENTAL REPORT / REGISTERED MODEL, PROBES AND THRESHOLDS}
{\fontsize{28}{31}\selectfont Helike as a Probe of a\\Quantum Yang--Mills Model\par}
\vspace{6pt}
{\large A geometric response measured on the GPU, with its convergence limits retained.}

\textbf{Abstract.} A gauge-invariant probe derived from Helike's homogeneous
geometry is evaluated in a published SU(2) Yang--Mills matrix model on the round
three-sphere. An unrestricted colour-singlet calculation within each finite
basis supplies the reference spectrum. At the largest registered cutoff,
the directional channel places 99.154\% of its excited response in the lowest
computed level at $g^2=2$. At $g^2=1$ and $0.5$, a separately registered scalar
composite reaches the lowest level instead. Fifteen cases pass finite-matrix
validation and evidence replay. Cutoff drift and exterior residuals remain too
large for a converged model-level conclusion.

\fig{01-lowest-level-response}{Response carried by the lowest complete excited
level at $N=12$ (1,134 physical basis states). Each channel is normalized by its
own finite-basis excited norm after vacuum subtraction. Tiny opposite-channel
weights are below $10^{-25}$ as fractions; they are not asserted to be exact zeros.
All three cases remain \textbf{UNRESOLVED} under the convergence contract.}

\textbf{What this establishes.} The mapped observables produce a reproducible,
shape-selective response in an external quantum model. They are compared against
the whole physical finite spectrum, not a spectrum restricted to look like the probe.

\textbf{What remains open.} The identity and energy of the true lowest excitation
have not settled under basis growth. The matched-control efficiency comparison
was therefore not run. This report makes no claim of a continuum mass gap,
computational advantage, or a result for the full four-dimensional theory.

\newpage
\pagetitle{01 / EXACT GEOMETRY}{The probe preserves a direction, not just a number}
The model coordinate $M\in\mathbb{R}^{3\times3}$ has spatial rows and colour
columns. Colour rotations act on the right. Thus $S=MM^T$ is gauge invariant.
The two real quadratures of the homogeneous member of Helike are
\[
 B_c=\frac1{\sqrt2}\begin{pmatrix}0&0&0\\0&1&0\\0&0&-1\end{pmatrix},\qquad
 B_s=\frac1{\sqrt2}\begin{pmatrix}0&0&0\\0&0&1\\0&1&0\end{pmatrix}.
\]
\fig{02-probe-geometry}{Exact directional signatures of the two declared tensors.
Colour is $n^TBn$ on $\|n\|=1$, with the same camera and colour scale in both
panels. These are calculated quadratic forms, not simulated vacuum shapes or
quantum eigenfunctions. The sphere's radius is constant; only its colour changes.
Coordinates use the registered spatial order $(0,1,2)$.}

Define $O_a(M)=\operatorname{tr}(B_aMM^T)$. At the mapped classical reference
connection $M_*=I/2$,
\[
 \nabla_M O_a\big|_{M_*}=2B_aM_*=B_a.
\]
This is the exact geometric link: the observable's first response is the original
direction. $M_*$ is a classical saddle used for calibration, not the assumed
quantum vacuum. The latter is independently solved across the declared domain.

Rotating these tensors spans all five symmetric trace-free spatial directions.
Their observables $O_A$ form the \textbf{directional spin-2 channel, H2}.
A separate scalar channel is their squared strength:
\[
 O_0=\sum_{A=1}^5 O_A^2
     =\operatorname{tr}(S^2)-\tfrac13(\operatorname{tr}S)^2.
\]
\textbf{H0 is a quartic composite, not the original linear eigenfield.}
From a spin-0 vacuum, H2 couples to spin 2 and H0 to spin 0. Neither observable
is expected to detect every symmetry sector. Only the homogeneous classical
member is translated here; the Berger floor theorem is not imported as a
quantum conclusion on the round sphere.

\newpage
\pagetitle{02 / MEASURED RESPONSE}{Where the geometry lands in the spectrum}
For the computed vacuum $\Omega$, let $Q$ remove its component and let $P_n$
project onto a complete excited energy level. The measured response is
\[
 \phi_A=QO_A\Omega,\qquad
 W_n=\frac{\sum_A\|P_n\phi_A\|^2}{\sum_A\|\phi_A\|^2}.
\]
Degenerate levels are scored as whole subspaces, not by choosing one eigenvector.
The full complex transition kernels are retained in the evidence.

\fig{03-spectral-response}{Finite-basis spectral fingerprints at $N=12$.
All complete levels with excitation energy at most $10/R$ are shown; weights
are normalized over \emph{all} finite levels, including those beyond the plotted
window. H2 (top) is concentrated in its first accessible spin-2 level at all
three couplings. H0 (bottom) spreads over several scalar levels. Dashed lines
mark the independently computed absolute gap, which need not be the probe's
first accessible gap.}

\begin{center}\small
\setlength{\tabcolsep}{5pt}
\begin{tabular}{@{}rrrrrr@{}}\toprule
$g^2$ & $E_0R$ & $\Delta_{\rm abs}R$ & $J_{\rm lowest}$ & $\Delta_{\Htwo}R$ & $\Delta_{\Hzero}R$\\\midrule
1 & 5.433212201 & 2.289814624 & 0 & 2.642429574 & 2.289814624 \\
0.5 & 4.891790030 & 1.566333799 & 0 & 2.339593687 & 1.566333799 \\
2 & 6.259094501 & 3.094157392 & 2 & 3.094157392 & 3.153908889 \\
\bottomrule\end{tabular}

\end{center}
At $g^2=1$ and $0.5$, the strong directional response is \emph{above} the absolute
lowest level. At $g^2=2$, it coincides with the lowest five-state level at this
cutoff. ``Accessible'' uses the registered weight threshold $10^{-4}$; smaller
weights remain in the evidence. Nine displayed decimals describe the finite
matrix computation, not nine-digit accuracy for the untruncated model.

\newpage
\pagetitle{03 / REFINEMENT HISTORY}{The answer must survive a larger space of states}
The reference basis includes every colour-invariant polynomial through degree
$N$, multiplied by a Gaussian. It retains all spatial spin sectors and the
determinant-odd states. Increasing $N$ adds possibilities; it does not steer
the solver toward Helike.

\fig{04-cutoff-history}{Registered cutoff history. (A) Absolute gaps at all
three couplings. (B) The first robustly accessible scalar and directional levels
at the strongest coupling. At $N=8$ the scalar is lower, whereas the directional
level is lower at $N=10$ and $12$. Lines connect computed cutoffs only; they are
neither interpolated eigenstates nor an extrapolation to infinite cutoff.}

\begin{center}\small
\setlength{\tabcolsep}{6pt}
\begin{tabular}{@{}rrrrrr@{}}\toprule
$N$ & Dimension & $\Delta R$, $g^2=1$ & $\Delta R$, $g^2=0.5$ & $\Delta R$, $g^2=2$ & $J$, $g^2=2$\\\midrule
4 & 29 & 2.867846746 & 2.037058771 & 4.351929591 & 2 \\
6 & 91 & 2.465167225 & 1.774780654 & 3.603696384 & 2 \\
8 & 238 & 2.379861209 & 1.667228247 & 3.264328103 & 0 \\
10 & 546 & 2.319649058 & 1.593988390 & 3.179460361 & 2 \\
12 & 1,134 & 2.289814624 & 1.566333799 & 3.094157392 & 2 \\
\bottomrule\end{tabular}

\end{center}

\textbf{The strongest-looking result has a necessary qualification.}
The 99.154\% directional weight at $N=12$ cannot be detached from the change in
level ordering at $N=8$. This is unresolved truncation dependence, not evidence
of a physical phase transition.

The registration requires two successive increments for which $E_0$, $E_1$
and each quoted gap change by no more than $10^{-3}\max(1,|\text{quantity}|)$,
and channel weights by no more than $0.01$. The survey does not meet these rules.
A difference of variational energy estimates is also not automatically an upper
or lower bound on the true gap.

\textbf{Boundary of this run.} $N=12$ was the only authorized extension beyond
$4,6,8,10$. The survey stopped there. No coupling, probe, Hamiltonian or acceptance
threshold was changed after inspecting the results.

\newpage
\pagetitle{04 / ACCEPTANCE}{Accurate finite matrices; unresolved model convergence}
Finite-matrix agreement answers whether the included problem was solved
accurately. Exterior checks ask whether too much of the full operator's action
still falls outside that included space. These are different questions.

\fig{05-convergence-gates}{Final-cutoff diagnostics divided by their fixed
acceptance limits. Ratios at or below the dashed line pass that individual test.
(A) Worst full Hamiltonian residual among the vacuum, lowest level and first
accessible probe levels, divided by $0.01\max(1,\Delta_{\rm abs})$.
(B) Discarded squared probe response divided by $0.001$.
The shaded region is the passing range. Every coupling fails the residual test;
the scalar tail fails in all three.}

\begin{center}\small
\begin{tabular}{@{}p{3.1in}p{3.05in}@{}}\toprule
\textbf{Verified in this experiment} & \textbf{Not established by this experiment}\\\midrule
20 focused tests; 15 exact symbolic identities & A cutoff-converged model gap\\
8 mechanism-removal checks rejected defects & The lowest level's eventual ordering\\
Independent physical-space construction at $N=4$ & A Helike efficiency advantage\\
CPU/GPU spectra and low-level projectors agree & A cross-model universality result\\
Matrix-exponential and spectral kernels agree & A four-dimensional continuum mass gap\\
All 15 cases replay from stored evidence & All HP00--HP10 research gates closed\\\bottomrule
\end{tabular}
\end{center}

The maximum scaled CPU/GPU spectral difference was $1.161\times10^{-15}$;
the maximum low-level projector difference was $1.646\times10^{-12}$, both below
the registered $10^{-10}$ limit. This high numerical precision does not erase
the exterior residuals shown above.

\textbf{Verdict: UNRESOLVED, with reproducible finite-cutoff observations.}
The random-seed and scalar-trace utility comparisons remain \texttt{NOT\_RUN}
because their convergence prerequisite failed. There is no hidden failed
comparator from which an advantage is inferred. A new approximation strategy
or a larger run requires a separate dated amendment.

\newpage
\pagetitle{05 / METHODS AND EVIDENCE}{A published Hamiltonian and a replayable calculation}
The external model is the pure SU(2) sphere reduction in Ref.~[1], with the
radius convention checked against Ref.~[2]. On the registered regular Cartesian
colour-singlet domain, setting $q=M/g$ gives
\[
 H=\frac1R\left[-\frac12\Delta_q+\frac12\operatorname{tr}(qq^T)
 -3g\det q+\frac{g^2}{4}
 \left\{(\operatorname{tr}(qq^T))^2-\operatorname{tr}((qq^T)^2)\right\}\right].
\]
The runs use $R=\hbar=1$, positive $g$, no added mass or fitted energy offset,
and no extra boundary extension. The two classical wells remain in the
Cartesian domain. The signed cubic determinant is not replaced by its absolute value.

\textbf{Basis and independent reference.}
The physical polynomial basis is generated by the six entries of $qq^T$ and
$\det q$, with determinant powers zero or one. The special orthogonal invariant
theorem and $(\det q)^2=\det(qq^T)$ justify this spanning construction [3].
Exact Wick ordering is a triangular change of this basis, not a dropped direction.
The overlap matrix is retained in the generalized eigenproblem. An independent
$N=4$ oscillator construction imposes the three Gauss generators and assembles
the Hamiltonian before projection. Direct quartic projection is used for H0;
squaring already projected quadratic probes would give a different operator.

\textbf{Computation.}
Double-precision CUDA solves and matrix exponentials ran on an NVIDIA GeForce
RTX 5070 Laptop GPU, with independent CPU comparisons. CUDA requests could not
fall back to CPU. The worker survey took 72.817 seconds including construction
checks, excluding the separate evidence replay. Peak sampled host memory was
2.4165 GiB; peak PyTorch GPU allocation was 157.992 MiB and reservation 226 MiB.
These device figures exclude driver/context allocations and other applications.
They are resource measurements, not a speedup benchmark.

\textbf{Provenance.}
Run \runid{} retains matrices, normalized eigenvectors, metric, angular
Casimir, observables, complete-level kernels, cutoffs, residuals and source hashes.
The source snapshot is included, because the computation predates the publishing
commit. Evidence replay returned \texttt{VERIFIED}, 15 cases,
\texttt{scientific\_acceptance: false}. The base source revision is:\\
{\small\texttt{ed5d3cf5209bb8cecd6d97c4f0df3065af90363e}}. Per-file hashes identify
the actual experimental implementation. No server endpoint was changed or deployed.

\textbf{Reuse.}
Every figure is supplied as PDF, SVG and PNG. The captions and equations are
editable LaTeX. CSV and JSON exports retain the plotted values, including tiny
weights. The accompanying README gives one-command figure/report generation
and the evidence-replay command. Numerical tables are generated from the receipt.

\vspace{5pt}
{\small\raggedright
\textbf{References}\par
[1] A.~P. Balachandran, S. Vaidya and A.~R. de Queiroz,
\emph{A Matrix Model for QCD}, arXiv:1412.7900v1,
\href{https://doi.org/10.1142/S0217732315500807}{doi:10.1142/S0217732315500807}.
\par [2] N. Acharyya and A.~P. Balachandran,
\emph{Matrix Model of QCD: Edge Localized Glue Balls and Phase Transitions},
\href{https://arxiv.org/abs/1702.06430v2}{arXiv:1702.06430v2}, Eq.~2.14.
\par [3] H. Kraft and C. Procesi, \emph{Classical Invariant Theory},
\href{https://dmi.unibas.ch/fileadmin/user_upload/dmi/Personen/Kraft_Hanspeter/Classical_Invariant_Theory.pdf}{Theorem 10.2}.
}

\newpage
\pagetitle{06 / NEXT EXPERIMENTAL LAYER}{Use GIGI to connect evidence across models}
\textbf{Plan, not a deployed capability.}
The model supplies observations. GIGI can organize those observations into
versioned bundles, compare compatible geometric signatures and expose which
features survive refinement. The goal is to follow a physical pattern through
changing approximations, not merely line up similarly sized gap numbers.

\begin{center}\small
\begin{tabular}{@{}p{1.35in}p{4.85in}@{}}\toprule
\textbf{Evidence layer} & \textbf{What must be retained}\\\midrule
Model contract & Hamiltonian, domain, gauge group, geometry, units, boundary conditions and normalization.\\[4pt]
Run and cutoff & Coupling, basis/mesh size, source and array hashes, resource data, status and all failed gates.\\[4pt]
Complete level & Energy, multiplicity, symmetry sector, projector identity and residual; not an arbitrary eigenvector label.\\[4pt]
Probe response & Full transition kernel, normalized spectral weights, vacuum-subtracted norm and discarded response.\\[4pt]
Comparison map & The explicit bridge between models or cutoffs, its preserved quantities, errors and validity domain.\\\bottomrule
\end{tabular}
\end{center}

\textbf{First deliverable: the present 15 cases.}
Import the verified receipt and immutable array references into isolated GIGI
bundles. Recover every gap and weight exactly. Follow the scalar and tensor
branches separately, so the $N=8$ ordering change cannot be hidden by sorting
each run and calling its first row the same state.

\textbf{Second deliverable: a measured geometric through-line.}
Within a common configuration space, embed successive cutoff subspaces into a
common metric and compare complete-level projectors. Across genuinely different
models, require a declared map or compare a common set of physical observables.
The shared signature can include dimensionless energies, spin labels, spectral
weight distributions and correlation kernels at fixed dimensionless times.
An apparent match must survive normalization, uncertainty and adversarial controls.

\textbf{Third deliverable: refinement before universality.}
Test convergence within each model separately. Then test whether the retained
patterns agree across compatible models, on held-out cases. Raising this
polynomial cutoff approaches the selected matrix model; it does not restore
the spatial field modes omitted by that model. A continuum claim requires its
own refinement parameters and limits.

The detailed plan is in
\texttt{docs/research/ym-evidence-continuum-gigi-plan-2026-11-23.md}.
It separates current source-backed GIGI operations from the missing adapters
and scientific gates. No graph spectral gap, PCA variance, or generic database
curvature is relabelled as a Yang--Mills excitation energy.

\textbf{Success criterion.}
GIGI should let a researcher ask: \emph{Is this the same geometric excitation
surviving a change of approximation, or only a similar number?} The answer must
identify the retained pattern, its comparison map and the evidence that could
falsify the match. It may also correctly return \emph{not yet comparable}.
\end{document}
